\documentclass[11pt,a4paper]{article}

% Packages
\usepackage[utf8]{inputenc}
\usepackage[T1]{fontenc}
\usepackage[margin=0.75in]{geometry}
\usepackage{enumitem}
\usepackage{hyperref}
\usepackage{titlesec}
\usepackage{xcolor}

% Formatting
\pagestyle{empty}
\setlength{\parindent}{0pt}
\setlist[itemize]{leftmargin=*, noitemsep, topsep=3pt}

% Section formatting
\titleformat{\section}{\Large\bfseries}{}{0em}{}[\titlerule]
\titlespacing*{\section}{0pt}{12pt}{6pt}

\titleformat{\subsection}[runin]{\bfseries}{}{0em}{}
\titlespacing*{\subsection}{0pt}{8pt}{0pt}

% Colors
\definecolor{linkcolor}{RGB}{0,102,204}
\hypersetup{
    colorlinks=true,
    urlcolor=linkcolor
}

\begin{document}

% Header
\begin{center}
    {\LARGE \textbf{ERIC MAYEUX}}\\
    \vspace{4pt}
    {\large Gestionnaire de Projet -- Projets Stratégiques | Strategic Project Manager}\\
    \vspace{4pt}
    Montréal, QC, Canada\\
    \href{mailto:mayeux.e@gmail.com}{mayeux.e@gmail.com} | 438-884-7645
\end{center}

\vspace{8pt}

\section*{Profil Professionnel}
\noindent
\textbf{Gestionnaire de projet stratégique} avec 10+ ans d'expérience en \textbf{coordination d'initiatives complexes} et \textbf{gouvernance multi-parties prenantes} dans le secteur bancaire. Expert en \textbf{gestion des risques}, suivi de performance et \textbf{ateliers de coopération} avec des équipes fonctionnelles variées. Maîtrise des méthodes de gestion de projet (Agile, hybride) et des outils PMO. Certifié \textbf{PMP} (en cours), CSM. Bilinguisme français/anglais.

\section*{Compétences Techniques}

\textbf{Gestion de projets stratégiques:} Coordination d'initiatives stratégiques, portée, budget, échéanciers, chemin critique, jalons, livrables, gestion matricielle

\textbf{Gouvernance \& Parties prenantes:} Gouvernance multi-parties prenantes (internes/externes), ateliers de coopération, comités exécutifs, rapports d'avancement

\textbf{Gestion des risques \& enjeux:} Identification, documentation, plans d'atténuation, suivi, post-mortem, leçons apprises

\textbf{Outils PMO:} Jira, MS Project, SharePoint, Confluence, Datadog, Splunk, tableaux de bord KPI

\textbf{Méthodologies:} Agile, Scrum, SAFe, Kanban, Waterfall, hybride, PMP

\textbf{Intelligence Artificielle:} Agents IA, automatisation, RPA

\textbf{Architecture \& Cloud:} AWS Cloud Practitioner, Architectures Distribuées

\vspace{6pt}

\section*{Réalisations}
\begin{itemize}
    \item \textbf{Coordination d'initiatives stratégiques} multi-équipes à la BNC : planification, exécution, livraison
    \item Animation d'\textbf{ateliers de coopération} pour faire avancer les activités de planification complexes
    \item Mise en place d'un suivi rigoureux : \textbf{chemin critique}, plans d'atténuation, jalons
    \item \textbf{Gouvernance} avec multiple parties prenantes : comités, rapports, escalades
    \item Suivi du \textbf{post-mortem} et partage des leçons apprises à la clôture de chaque projet
    \item Déploiement d'\textbf{agents IA} pour l'automatisation des tâches de gestion de projet
\end{itemize}

\section*{Expérience Professionnelle}

\subsection*{Scrum Master -- Gestion de projet | Project Manager} \hfill \textit{Novembre 2025 -- Présent}\\
\textbf{Banque Nationale du Canada (BNC)} -- Montréal, QC\\
Pilote des \textbf{projets stratégiques} TI : coordination, gouvernance, suivi de performance.
\begin{itemize}
    \item Coordonne les \textbf{ateliers de coopération} pour faire progresser les activités de planification
    \item Assure l'exécution, le suivi et la \textbf{livraison} des projets (portée, budget, échéanciers, risques)
    \item Anime des séances de travail avec les parties prenantes pour \textbf{résoudre les problèmes}
    \item Développe un échéancier supportant les \textbf{jalons} ; détermine le \textbf{chemin critique}
    \item Participe à la mise en œuvre du processus de \textbf{gestion des risques} et favorise une culture de saine gestion
    \item Assure le \textbf{post-mortem} et partage les leçons apprises avec les parties concernées
\end{itemize}

\subsection*{Intégrateur Sénior Qualité | Senior Quality Integrator} \hfill \textit{Octobre 2023 -- Novembre 2025}\\
\textbf{Banque Nationale du Canada (BNC)} -- Montréal, QC\\
\begin{itemize}
    \item \textbf{Gouvernance} des pratiques qualité pan-train ; coordination multi-sites (Thaïlande)
    \item Tableaux de bord et \textbf{rapports d'avancement} pour la direction et les parties prenantes
    \item Animation d'\textbf{ateliers} et formations ; gestion des risques qualité
\end{itemize}

\subsection*{Intégrateur Qualité -- SDET | Quality Integrator -- Software Development Engineer in Test} \hfill \textit{Janvier 2022 -- Octobre 2023}\\
\textbf{Banque Nationale du Canada (BNC)} via Groupe Astek -- Montréal, QC\\
\begin{itemize}
    \item \textbf{Gestion des risques} qualité ; métriques, pyramide de test, sécurité, stabilité
    \item Stack : Selenium, AppliTools, RestAssured, JMeter, Gatling, K6
\end{itemize}

\subsection*{QA Automaticien | QA Automation Engineer} \hfill \textit{Juin 2021 -- Décembre 2021}\\
\textbf{AODocs} -- Paris, France\\
\begin{itemize}
    \item Automatisation tests non-régression ; POC Flutter ; rapports Allure
\end{itemize}

\subsection*{Automaticien CI/CD | DevOps Automation Engineer} \hfill \textit{Janvier 2020 -- Juin 2021}\\
\textbf{ENGIE DIGITAL} via Groupe Hénix -- Paris, France\\
\begin{itemize}
    \item Création CI/CD ; pilotage projets Jira, KPIs ; \textbf{coordination} d'équipes
\end{itemize}

\subsection*{Product Owner} \hfill \textit{Juin 2019 -- Septembre 2019}\\
\textbf{Hénix} -- Montrouge, France\\
\begin{itemize}
    \item Coordination \textbf{stratégique} Jira/Squash (8 500 utilisateurs) ; gestion du backlog
\end{itemize}

\subsection*{Développeur Logiciel | Software Developer} \hfill \textit{Mai 2017 -- Mai 2019}\\
\textbf{MEDIAPOST} via Groupe Hénix -- Paris, France\\
\begin{itemize}
    \item Développement back office, webservices, reporting Excel
\end{itemize}

\subsection*{Consultant QA \& Automatisation | QA Automation Consultant} \hfill \textit{Avril 2016 -- Avril 2017}\\
\textbf{ENGIE IT} via Groupe Hénix -- Saint-Ouen, France\\
\begin{itemize}
    \item Pilotage opérationnel, gestion des incidents, administration HP ALM
\end{itemize}

\section*{Certifications \& Formations}

\textbf{PMP Certification} (en cours) \hfill \textit{2026}

Project Management Professional -- Project Management Institute (PMI)

\textbf{AWS Cloud Practitioner} \hfill \textit{2026}

Amazon Web Services (AWS)

\textbf{Certified Scrum Master (CSM)} \hfill \textit{2024}

Scrum Alliance

\textbf{Jira ACP-600 -- Project Administration in Jira Server} \hfill \textit{2019}

Atlassian Certified Professional

\textbf{Cursus Automaticien et DevOps} \hfill \textit{2019}\\
\textit{AFCEPF -- Montrouge, France}

\textbf{Cursus Architecture logiciel \& Analyste informaticien - Certifié Master II} \hfill \textit{2017}\\
\textit{AFCEPF-HENIX -- Malakoff, France}

\section*{Formation Académique}

\textbf{Licence en Gestion (Bachelor's Degree in Management)} \hfill \textit{2011}\\
École Supérieure de Gestion (E.S.G) -- Paris, France

\section*{Compétences Linguistiques}

\textbf{Français:} Langue maternelle (Native)\\
\textbf{Anglais:} Niveau Avancé (Advanced / Professional Working Proficiency)

\end{document}
